% 19-memory.tex — Clause 19: Memory model
\chapter{Memory model}
\label{sec:19-memory}
\index{memory model}

\section{Overview}
\label{sec:mem.overview}

\pnum
Lain's memory model is founded on three invariants:
\begin{enumerate}
  \item \textbf{Deterministic allocation} — every allocation is explicit
        and predictable; the implementation never implicitly allocates
        heap memory.
  \item \textbf{Zero runtime overhead} — there is no garbage collector
        and no reference-counting mechanism.
  \item \textbf{C99 layout compatibility} — the physical layout of all
        objects follows ISO/IEC~9899 (C17) rules exactly.
\end{enumerate}

\pnum
This clause specifies the abstract notions of \textit{object},
\textit{storage duration}, \textit{object representation}, and
\textit{address}.  The abstract machine described here is realised by any
conforming translation output (\S\,\ref{sec:conformance}).

\section{Objects}
\label{sec:mem.objects}
\index{object}

\pnum
An \textit{object} is a contiguous region of storage that holds a value of
a declared type.  Every variable declaration introduces exactly one object.
Fields of a struct introduce sub-objects.  Elements of a fixed-size array
introduce element sub-objects.

\pnum
Every object has:
\begin{itemize}
  \item a \textit{declared type} (\S\,\ref{sec:07-types});
  \item an \textit{ownership mode} — shared, \lain{var}, or \lain{mov}
        (\S\,\ref{sec:own.modes});
  \item a \textit{linear state} — \textsc{Uninit}, \textsc{Unconsumed},
        or \textsc{Consumed} (\S\,\ref{sec:own.states});
  \item a \textit{storage duration} (\S\,\ref{sec:mem.storage});
  \item an \textit{address} — the byte address of its first storage unit,
        aligned as required by the target platform.
\end{itemize}

\section{Storage duration}
\label{sec:mem.storage}
\index{storage duration}

\pnum
Three storage durations are defined:

\begin{description}[leftmargin=2em]
  \item[Automatic duration]
    Objects declared inside a function or block body (including function
    parameters) have automatic storage duration.  Storage is allocated when
    the enclosing scope is entered and released when that scope exits.

  \item[Static duration]
    Objects declared at module scope have static storage duration.  Their
    storage persists for the entire program execution.
    \idbehav{The value of a static-duration object whose declaration
    lacks an initializer is implementation-defined; the reference
    implementation zero-initializes such objects via C99 static
    initialization rules}.

  \item[Heap (explicit) duration]
    Objects allocated through C interop (\lain{malloc} or equivalent)
    have programmer-controlled duration.  The ownership system
    (\S\,\ref{sec:11-ownership}) enforces that \lain{mov} pointers to
    heap memory are consumed exactly once~(\diagcode{E003}).
\end{description}

\begin{constraint}
  A \lain{return} statement shall not return a reference whose referent
  has automatic storage duration~(\diagcode{E010}).
\end{constraint}

\section{Stack allocation}
\label{sec:mem.stack}
\index{stack allocation}

\pnum
All local variables, struct values, and fixed-size arrays are allocated on
the call stack.  No implicit heap allocation is ever performed by the
compiler.

\begin{laincode}
var x int = 42     // stack: sizeof(int32_t) bytes
var arr int[100]    // stack: 100 * sizeof(int32_t) bytes
var p Point         // stack: sizeof(Point) bytes
\end{laincode}

\pnum
Each function call produces a \textit{stack frame} containing:
\begin{itemize}
  \item all local variables declared in the function;
  \item all function parameters (by value or by pointer, per the calling
        convention of \S\,\ref{sec:fn.distinction});
  \item compiler-generated temporaries required for expression evaluation.
\end{itemize}

\pnum
The stack frame is automatically reclaimed when the function returns.
\idbehav{The implementation's behaviour when the call stack is exhausted
is implementation-defined; the reference implementation inherits the
platform's stack overflow signal}.

\section{Heap allocation}
\label{sec:mem.heap}
\index{heap allocation}

\pnum
Heap allocation is performed exclusively through C interop
(\S\,\ref{sec:17-interop}).  The idiomatic pattern declares the C
\texttt{malloc}/\texttt{free} pair as \lain{extern proc}:

\begin{laincode}
c_include "<stdlib.h>"
extern proc malloc(size usize) mov *void
extern proc free(ptr mov *void)
\end{laincode}

\pnum
Because \lain{malloc} returns \lain{mov *void}, the caller receives sole
ownership of the allocated block.  Because \lain{free} consumes
\lain{mov *void}, calling \lain{free} transfers that ownership to the
allocator.  The ownership system therefore enforces statically:
\begin{itemize}
  \item \textbf{No use-after-free}: using the pointer after \lain{free}
        is a use-after-move~(\diagcode{E001}).
  \item \textbf{No double-free}: a second call to \lain{free} on the
        same pointer is a double-consume~(\diagcode{E002}).
  \item \textbf{No leak}: failing to call \lain{free} on a \lain{mov}
        pointer is an unconsumed linear value~(\diagcode{E003}).
\end{itemize}

\section{Value semantics}
\label{sec:mem.value-semantics}
\index{value semantics}

\pnum
The behaviour of assignment and function argument passing is determined
by the type's linearity and the ownership mode:

\begin{longtable}{@{}lll@{}}
\toprule
\textbf{Type} & \textbf{Assignment \texttt{b = a}} & \textbf{Notes} \\
\midrule
\endfirsthead
\toprule
\textbf{Type} & \textbf{Assignment \texttt{b = a}} & \textbf{Notes} \\
\midrule
\endhead
\bottomrule
\endlastfoot
Primitives (\lain{int}, \lain{bool}, \ldots) & Bitwise copy & Always allowed \\
Non-linear struct                             & Field-by-field copy & Allowed \\
Linear struct (has \lain{mov} field)          & Forbidden & Requires \lain{mov} \\
Non-linear array                              & Element-wise copy & Allowed \\
Linear array                                  & Forbidden & Requires \lain{mov} \\
Slice (\lain{[T]})                            & Copy \texttt{\{data, len\}} pair & Always allowed \\
Raw pointer (\lain{*T})                       & Copy address & Always allowed \\
\end{longtable}

\begin{constraint}
  Copying a linear type without an explicit \lain{mov} expression is
  \illformed{}~(\diagcode{E006}).
\end{constraint}

\section{Object representation}
\label{sec:mem.representation}
\index{object representation}

\pnum
The \textit{object representation} of a Lain type is its C99 equivalent
as defined in \S\,\ref{sec:interop.types}.  Every object is a sequence of
\lain{u8} bytes whose count equals the C99 \texttt{sizeof} of the
corresponding C type.

\pnum
\textit{Value equality} is defined per type:
\begin{itemize}
  \item Primitives: bitwise equality of the value representation.
  \item Structs: field-by-field value equality.
  \item Slices: equality of both the data pointer and the length.
  \item Raw pointers: address equality.
  \item ADTs: tag equality and, when tags match, field-by-field equality
        of the active variant.
\end{itemize}

\section{Alignment and layout}
\label{sec:mem.alignment}
\index{alignment}
\index{struct layout}

\pnum
Lain inherits ISO/IEC~9899 (C17) alignment and padding rules:
\begin{itemize}
  \item Struct fields are laid out in declaration order.
  \item Each field is aligned to its natural alignment (the alignment of
        its C99 equivalent type).
  \item The compiler inserts padding between fields as necessary.
  \item The total size of a struct is rounded up to a multiple of its
        largest field alignment.
\end{itemize}

\pnum
Lain does not provide layout control annotations.  There is no
\texttt{\#pragma pack} equivalent, no \texttt{alignas}, and no
\texttt{repr(packed)}.  The layout is always the platform-default C99
layout.

\idbehav{The exact size and alignment of each type is implementation-defined,
following the C99 ABI of the target platform}.

\section{ADT and enum layout}
\label{sec:mem.adt-layout}
\index{ADT layout}
\index{enum layout}

\pnum
An enum type is emitted as a C \texttt{enum} wrapped in a named struct so
that each enum type is nominally distinct in the C translation:

\begin{verbatim}
typedef enum { Color_Red, Color_Green, Color_Blue } Color_Tag;
typedef struct { Color_Tag tag; } Color;
\end{verbatim}

\pnum
An ADT type uses a niche-based encoding whenever a zero-cost
representation is derivable. The compiler performs \textit{aggressive
multi-slot niche optimization} (D-Niche): it explores the full
sentinel space of the payload and packs as many empty variants as
fit, without any tag byte. When the pool is insufficient to cover
every empty variant, the layout falls back to a tagged union and the
compiler emits \diagcode{W120} (best-effort warning; not an error).

\begin{verbatim}
// Tagged-union fallback (only when no niche is derivable):
typedef struct {
    enum { Shape_Circle_tag, Shape_Rect_tag } tag;
    union {
        struct { float r; } Circle;
        struct { float w; float h; } Rect;
    };
} Shape;
\end{verbatim}

\pnum
\textit{Niche encoding (zero-cost).}  When the payload of one variant
contains a bit pattern that cannot occur in valid values---for instance,
\texttt{0x0} for a pointer field, or a value outside a VRA-proven
range for an integer field---that pattern is reserved to encode an
alternative variant.  No explicit tag is emitted.

\begin{verbatim}
// Niche encoding for Option(*File):  size = sizeof(*File)
typedef File *Option_Ptr_File;
//  Some(p)   -> non-null pointer p
//  None      -> 0x0
\end{verbatim}

\pnum
The sentinel pool size depends on the target. For pointer payloads,
the pool is the zero-page-unmappable region of the target's address
space, stepped by the pointer alignment (e.g.\ 16384 slots for
x86\_64-linux with alignment~8). For integer payloads, the pool is
the type's full range minus any VRA-narrowed sub-range:
\lain{type Pct = i32 >= 0 and <= 100} as a payload yields a pool
of $\approx 4.3 \times 10^9$ slots (every integer outside $[0,100]$).
For \lain{bool}, the pool is the 254 byte patterns outside $\{0, 1\}$.

\pnum
Niche allocation cascades through nested enum payloads. If
\lain{Option(*T)} reserves the value~\texttt{0} for \lain{None}, a
surrounding \lain{Result(Option(*T), Err4)} sees the remaining slots
$(8, 16, 24, \ldots)$ and packs its own empty variants into them.
Layout of the entire result still occupies one pointer-sized word.

\pnum
\textit{Multi-payload niche.} An enum with exactly two payload variants
where one payload (the \textit{primary}) has a non-empty sentinel pool
and the other payload (the \textit{secondary}) has a single field whose
type is an \textit{enum of pure-empty variants} is encoded by mapping
each sub-variant into a sentinel of the primary's pool. For
\lain{Result(*T, FileErr)} with \lain{FileErr = enum \{ NotFound,
Permission, OOM \}}, the encoding on a host with pointer alignment~8 is
\lain{Ok(f)} = any valid pointer, \lain{Err(NotFound)} = sentinel~0,
\lain{Err(Permission)} = sentinel~8, \lain{Err(OOM)} = sentinel~16. The
secondary constructor emits a stride-multiplication; the case-match
discriminator emits a corresponding stride-division for binding-style
patterns.

\pnum
The compiler never refuses compilation to satisfy a layout
optimization: when the pool is insufficient, \diagcode{W120} is
emitted with suggestions (constrain the payload, change to a type
with larger sentinel space, or split the enum), and a tag byte is
added. The programmer decides whether to accept the overhead.

\pnum
The CLI flag \texttt{--target=<triple>} selects the target's sentinel
pool size (host auto-detection is the default). Supported triples
include \texttt{x86\_64-linux-gnu}, \texttt{aarch64-linux-gnu},
\texttt{x86\_64-windows-msvc}, and \texttt{cortex-m4-bare} (the
last has zero pool: no MMU means address~0 may be a valid vector
table entry).

\section{No garbage collector}
\label{sec:mem.no-gc}
\index{garbage collector!absence of}

\pnum
Lain has no garbage collector.  Consequently:
\begin{itemize}
  \item execution is never interrupted for memory management;
  \item deallocation occurs at statically deterministic program points;
  \item no reference counting or heap tracing is performed at runtime.
\end{itemize}

\pnum
This makes Lain suitable for real-time, embedded, and safety-critical
applications where garbage-collection pauses are unacceptable.

\section{Memory safety in the safe fragment}
\label{sec:mem.safety}
\index{memory safety}

\pnum
In the \textit{safe fragment} (code outside \lain{unsafe} blocks), the
following memory safety properties hold by construction:

\begin{longtable}{@{}ll@{}}
\toprule
\textbf{Property} & \textbf{Enforcement mechanism} \\
\midrule
\endfirsthead
\toprule
\textbf{Property} & \textbf{Enforcement mechanism} \\
\midrule
\endhead
\bottomrule
\endlastfoot
No use-after-free         & \diagcode{E001} (use-after-move) \\
No double-free            & \diagcode{E002} (double-consume) \\
No leak of linear values  & \diagcode{E003} (unconsumed at scope exit) \\
No dangling references    & \diagcode{E010} (ref to automatic-duration object) \\
No out-of-bounds access   & \diagcode{E014} (VRA proof failure) \\
No null-pointer deref     & No raw pointers in safe fragment \\
No data races             & No concurrency in current version \\
\end{longtable}

\begin{constraint}
  Inside \lain{unsafe} blocks, pointer operations may violate the
  above properties if the programmer's safety assertion is incorrect
  (\S\,\ref{sec:unsafe.ub}).
\end{constraint}

\section{Compiler-internal arena (informative)}
\label{sec:mem.arena}
\index{arena!compiler-internal}

\begin{note}
  The reference implementation uses a bump-pointer arena internally for
  all AST nodes, symbol table entries, and diagnostic objects.  This
  arena is an implementation detail; it is never exposed to user code.
  User programs may implement arena allocators of their own through C
  interop.
\end{note}
